\documentclass[11pt,a4paper]{article}
\usepackage{../common/td_style}

\begin{document}

\makeuniversityheader{Corrigé des Travaux Dirigés : Série 03}{Corrélation \& Régression Linéaire en Biochimie}{\textcolor{BioTeal}{\textbf{CORRIGÉ DÉTAILLÉ}}}

\begin{solutionbox}{Exercice 1 : Dosage Protéique de Bradford \& Coefficient de Pearson}
\begin{enumerate}[label=\textbf{\alph*.}]
  \item \textbf{Sommes élémentaires ($n = 5$) :}
  \begin{itemize}
    \item $\sum x_i = 2.0 + 4.0 + 6.0 + 8.0 + 10.0 = \mathbf{30.0\,\mu\text{g/mL}}$ ($\bar{x} = 6.0$).
    \item $\sum y_i = 0.125 + 0.245 + 0.380 + 0.495 + 0.635 = \mathbf{1.880\,\text{OD}}$ ($\bar{y} = 0.376$).
    \item $\sum x_i^2 = 4 + 16 + 36 + 64 + 100 = \mathbf{220.0}$.
    \item $\sum y_i^2 = 0.015625 + 0.060025 + 0.144400 + 0.245025 + 0.403225 = \mathbf{0.86830}$.
    \item $\sum x_i y_i = (2 \times 0.125) + (4 \times 0.245) + (6 \times 0.380) + (8 \times 0.495) + (10 \times 0.635)$ \\
    $= 0.250 + 0.980 + 2.280 + 3.960 + 6.350 = \mathbf{13.820}$.
  \end{itemize}

  \item \textbf{Sommes des carrés corrigées :}
  \begin{itemize}
    \item $SS_{xx} = \sum x_i^2 - \frac{(\sum x_i)^2}{n} = 220.0 - \frac{30.0^2}{5} = 220.0 - 180.0 = \mathbf{40.00}$.
    \item $SS_{yy} = \sum y_i^2 - \frac{(\sum y_i)^2}{n} = 0.86830 - \frac{1.880^2}{5} = 0.86830 - 0.70688 = \mathbf{0.16142}$.
    \item $SP_{xy} = \sum x_i y_i - \frac{(\sum x_i)(\sum y_i)}{n} = 13.820 - \frac{30.0 \times 1.880}{5} = 13.820 - 11.280 = \mathbf{2.540}$.
  \end{itemize}

  \item \textbf{Coefficient de corrélation linéaire de Pearson ($r$) :}
  \begin{equation*}
    r = \frac{SP_{xy}}{\sqrt{SS_{xx} \cdot SS_{yy}}} = \frac{2.540}{\sqrt{40.0 \times 0.16142}} = \frac{2.540}{\sqrt{6.4568}} = \frac{2.540}{2.54102} = \mathbf{0.9996}
  \end{equation*}

  \item \textbf{Test de significativité de $r$ ($df = n - 2 = 3$) :}
  \begin{equation*}
    t_{\text{obs}} = \frac{r \sqrt{n - 2}}{\sqrt{1 - r^2}} = \frac{0.9996 \times \sqrt{3}}{\sqrt{1 - 0.9992}} = \frac{0.9996 \times 1.73205}{\sqrt{0.0008}} = \frac{1.7314}{0.02828} = \mathbf{61.22}
  \end{equation*}
  Puisque $t_{\text{obs}} = 61.22 \gg t_{\text{crit}}(3) = 5.841$, la corrélation linéaire entre la concentration en BSA et l'absorbance est \textbf{extrêmement hautement significative} ($p < 0.0001$).
\end{enumerate}
\end{solutionbox}

\begin{solutionbox}{Exercice 2 : Droite des Moindres Carrés \& Coefficient de Détermination $R^2$}
\begin{enumerate}[label=\textbf{\alph*.}]
  \item \textbf{Équation de la droite de régression $\hat{y} = ax + b$ :}
  \begin{itemize}
    \item Pente : $a = \frac{SP_{xy}}{SS_{xx}} = \frac{2.540}{40.00} = \mathbf{0.0635\,\text{OD}/(\mu\text{g/mL})}$.
    \item Ordonnée à l'origine : $b = \bar{y} - a \bar{x} = 0.376 - (0.0635 \times 6.0) = 0.376 - 0.381 = \mathbf{-0.005\,\text{OD}}$.
  \end{itemize}
  L'équation est : $\mathbf{\hat{y} = 0.0635 \, x - 0.005}$.

  \item \textbf{Signification biochimique :}
  \begin{itemize}
    \item Pente $a = 0.0635$ : C'est la **sensibilité analytique** de la méthode de Bradford. Chaque ajout de $1\,\mu\text{g/mL}$ de BSA induit une augmentation de $0.0635$ unité d'absorbance à $595\,\text{nm}$.
    \item Ordonnée $b = -0.005$ : Très proche de 0 (non significativement différente de 0), elle confirme l'excellence du zéro optique réalisé contre le blanc.
  \end{itemize}

  \item \textbf{Coefficient de détermination $R^2$ :}
  $R^2 = r^2 = (0.9996)^2 = \mathbf{0.9992}$ (soit $\mathbf{99.92\%}$).
  $99.92\%$ de la dispersion observée des absorbances est rigoureusement expliquée par les variations de concentration protéique.

  \item \textbf{Critère de validation qualité :}
  Puisque $R^2 = 0.9992 \ge 0.995$, la gamme étalon est formellement validée pour un usage en laboratoire d'analyse accrédité.
\end{enumerate}
\end{solutionbox}

\begin{solutionbox}{Exercice 3 : Dosage d'un Échantillon Inconnu \& Intervalle de Confiance}
\begin{enumerate}[label=\textbf{\alph*.}]
  \item \textbf{Concentration estimée dans l'échantillon dilué ($\hat{x}_0$) :}
  \begin{equation*}
    \hat{x}_0 = \frac{y_0 - b}{a} = \frac{0.410 - (-0.005)}{0.0635} = \frac{0.415}{0.0635} = \mathbf{6.535\,\mu\text{g/mL}}
  \end{equation*}

  \item \textbf{Erreur type d'estimation inverse ($s_{\hat{x}_0}$) :}
  \begin{align*}
    s_{\hat{x}_0} &= \frac{s_{y/x}}{a} \sqrt{1 + \frac{1}{n} + \frac{(y_0 - \bar{y})^2}{a^2 \cdot SS_{xx}}} \\
    &= \frac{0.00894}{0.0635} \sqrt{1 + \frac{1}{5} + \frac{(0.410 - 0.376)^2}{(0.0635)^2 \times 40.0}} \\
    &= 0.1408 \sqrt{1 + 0.20 + \frac{0.001156}{0.1613}} = 0.1408 \sqrt{1.2072} = 0.1408 \times 1.0987 = \mathbf{0.155\,\mu\text{g/mL}}
  \end{align*}

  \item \textbf{Intervalle de Confiance à 95\% :}
  \begin{align*}
    \text{IC}_{95\%} &= \hat{x}_0 \pm t_{0.05, \, 3} \times s_{\hat{x}_0} \\
    &= 6.535 \pm (3.182 \times 0.155) = 6.535 \pm 0.493 = \mathbf{[6.04 \,;\, 7.03]\,\mu\text{g/mL}}
  \end{align*}

  \item \textbf{Concentration tissulaire native :}
  En multipliant par le facteur de dilution 100 :
  $C_{\text{foie}} = 6.535 \times 100 = 653.5\,\mu\text{g/mL} = \mathbf{0.654\,\text{mg/mL}}$ (ou $\text{g/L}$).
\end{enumerate}
\end{solutionbox}

\begin{solutionbox}{Exercice 4 : Diagnostic des Résidus \& Limites du Modèle Linéaire}
\begin{enumerate}[label=\textbf{\alph*.}]
  \item \textbf{Définition et propriété des résidus :}
  $e_i = y_i - \hat{y}_i$ (écart entre l'absorbance réelle et l'absorbance prédite par la droite). Dans la régression par moindres carrés avec constante, la somme des résidus est obligatoirement nulle : $\sum e_i = 0$.

  \item \textbf{Forme en U inversé des résidus :}
  Elle révèle une **non-linéarité systématique**. En biochimie, au-delà d'une certaine concentration en protéines, les molécules de colorant libre Coomassie s'épuisent ou la lumière diffusée n'obéit plus à la loi de Beer-Lambert (phénomène de saturation optique). Le modèle linéaire n'est alors plus valide sur cette plage.

  \item \textbf{Résidus en entonnoir (Hétéroscédasticité) :}
  L'hypothèse d'homoscédasticité ($\sigma_e^2 = \text{constante}$) est violée : l'incertitude sur la mesure grandit avec la concentration. Les Moindres Carrés Ordinaires (MCO) accordent trop de poids aux fortes concentrations imprécises. La solution est la **Régression Linéaire Pondérée (WLS)** avec des poids inversement proportionnels à la variance : $w_i = 1 / s_i^2$ ou $w_i = 1 / x_i^2$.

  \item \textbf{Usage du coefficient de Spearman :}
  Le coefficient des rangs de Spearman doit être privilégié lorsque :
  \begin{itemize}
    \item La relation entre les variables est monotone mais non strictement linéaire (courbe sigmoïde ou hyperbolique).
    \item Les données comportent des valeurs aberrantes extrêmes ou ne suivent pas la loi normale.
  \end{itemize}
\end{enumerate}
\end{solutionbox}

\begin{solutionbox}{Exercice 5 : Limites Analytiques de Détection (LOD) et de Quantification (LOQ)}
\begin{enumerate}[label=\textbf{\alph*.}]
  \item \textbf{Calcul de la Limite de Détection (LOD) :}
  \begin{equation*}
    \text{LOD} = \frac{3.3 \cdot s_{y/x}}{a} = \frac{3.3 \times 0.00894}{0.0635} = \frac{0.02950}{0.0635} = \mathbf{0.465\,\mu\text{g/mL}}
  \end{equation*}

  \item \textbf{Calcul de la Limite de Quantification (LOQ) :}
  \begin{equation*}
    \text{LOQ} = \frac{10 \cdot s_{y/x}}{a} = \frac{10 \times 0.00894}{0.0635} = \frac{0.0894}{0.0635} = \mathbf{1.408\,\mu\text{g/mL}}
  \end{equation*}

  \item \textbf{Distinction biologique et réglementaire :}
  \begin{itemize}
    \item **LOD ($0.46\,\mu\text{g/mL}$) :** Plus faible concentration pouvant être distinguée avec certitude du bruit de fond du blanc ($p < 0.01$), mais sans précision chiffrée suffisante (CV $> 20\%$).
    \item **LOQ ($1.41\,\mu\text{g/mL}$) :** Plus faible concentration quantifiable avec une fidélité et une justesse acceptables selon les normes internationales (erreur relative $< 10\%$).
  \end{itemize}

  \item \textbf{Décision sur l'échantillon à $0.80\,\mu\text{g/mL}$ :}
  Puisque $\text{LOD} (0.46) < 0.80 < \text{LOQ} (1.41)$ :
  \textbf{Non !} On ne peut pas consigner une valeur numérique quantitative. La présence de protéines est formellement détectée, mais le rapport d'analyse doit porter la mention réglementaire : « \textit{Protéines détectées, inférieures à la limite de quantification ($< 1.41\,\mu\text{g/mL}$)} » ou bien le biologiste doit recommencer le dosage avec un échantillon moins dilué.
\end{enumerate}
\end{solutionbox}

\end{document}
